\documentclass{article}

% Language setting
% Replace `english' with e.g. `spanish' to change the document language
\usepackage[czech]{babel}
\usepackage{amsthm,thmtools,xcolor,amsmath,amssymb, mathtools}

% Set page size and margins
% Replace `letterpaper' with `a4paper' for UK/EU standard size
\usepackage[a4paper,top=2cm,bottom=2cm,left=3cm,right=3cm,marginparwidth=1.75cm]{geometry}

% Useful packages
\usepackage{enumerate}
\usepackage{graphicx}
\usepackage{tabularx}
\usepackage[colorlinks=true, allcolors=blue]{hyperref}
\graphicspath{images/}

%\usepackage{tikzit}
%\input{default.tikzstyles}

\title{DÚ Lineární algebra -- Sada 7}
\author{Jan Romanovský}

\begin{document}
\maketitle

\textbf{(7.2)} Když si explicitně zapíšeme kam zobrazí operátor obecný polynom $ax^2 + bx + c$:
$$\phi:ax^2+ bx + c \longmapsto 2(2ax+b)+x^2(2a)=2ax^2+4ax+2b$$

Zvolme si pro polynomy stupně nejvýše 2 kanonickou bázi $K = \left\{x^2, x, 1\right\}$. V této bázi zapíšeme operátor $\phi$ maticově jako:
$$\phi\left(\begin{pmatrix}
   a \\ b \\ c
\end{pmatrix}\right) = \begin{pmatrix}
   2a \\ 4a \\ 2b
\end{pmatrix}= \begin{pmatrix}
    2 & 0 & 0\\
    4 & 0 & 0\\
    0 & 2 & 0
\end{pmatrix}\begin{pmatrix}
   a \\ b \\ c
\end{pmatrix}.$$

Když najdeme vlastní čísla a vlastní vektory:
$$p(\lambda)=(2-\lambda)\lambda^2; \lambda_1 = 2, \lambda_2 = 0; M_1 = \text{span }\left\{\begin{pmatrix}
   1 \\ 2 \\ 2
\end{pmatrix}\right\}, M_2 = \text{span }\left\{\begin{pmatrix}
   0 \\ 0 \\ 1
\end{pmatrix}\right\}.$$

Mějme tedy $\textbf{v}_1^1 = (1, 2, 2)^T, \textbf{v}_1^2 = (0, 0, 1)^T$. Nyní budeme hledat řetízky:
\begin{align*}
    \textbf{v}_1^1 \xmapsto{f-2\text{id}} &\,\textbf{o}\\
    \textbf{v}_2^2 \xmapsto{\phantom{-2}f\phantom{\text{id}}} \textbf{v}_2^1 \xmapsto{\phantom{-2}f\phantom{\text{id}}} &\,\textbf{o}
\end{align*}

Přičemž nám stačí už najít jen poslední vektor $\textbf{v}_2^2$, který splňuje dvě následující podmínky:

$$\phi(\textbf{v}_2^2)=\textbf{v}_1^2,\, \textbf{v}_2^2 \in \text{Ker }\phi^2$$
$$\left(\begin{tabular}{ c c c | c }
    2 & 0 & 0 & 0\\
    4 & 0 & 0 & 0\\
    0 & 2 & 0 & 1
\end{tabular}\right)\implies\textbf{v}_2^2 \stackrel{?}{=} \begin{pmatrix}
   0 \\ \frac{1}{2} \\ 0
\end{pmatrix},$$
$$\text{Ker }\begin{pmatrix}
    2 & 0 & 0\\
    4 & 0 & 0\\
    0 & 2 & 0
\end{pmatrix}^2 = \text{span }\left\{\begin{pmatrix}
   0 \\ 1 \\ 0
\end{pmatrix}, \begin{pmatrix}
   0 \\ 0 \\ 1
\end{pmatrix}\right\},$$
$$\implies \textbf{v}_2^2 = \begin{pmatrix}
   0 \\ \frac{1}{2} \\ 0
\end{pmatrix}.$$

Našli jsme tedy bázi $B$ v níž je operátor v Jordanově tvaru:

$$B = \left\{\begin{pmatrix}
   1 \\ 2 \\ 2
\end{pmatrix}, \begin{pmatrix}
   0 \\ 0 \\ 1
\end{pmatrix}, \begin{pmatrix}
   0 \\ \frac{1}{2} \\ 0
\end{pmatrix}\right\}, \, [\phi]_B^B = \begin{pmatrix}
   2 & 0 & 0 \\
    0 & 0 & 1 \\
    0 & 0 & 0
\end{pmatrix}.$$
\end{document}
